\chapter{Actions, actors and local information}\label{ch:actors}
\question{Who acts, what is evaluated, and what information is allowed?}

The transfer has a source, a destination and a quantity. An actor can propose it. These are separate facts. A person may make several proposals; several people may cooperate on one process; an automatic device may perform the physical operation. The mathematical description of the action does not by itself assign ownership or responsibility.

For our simple world, write an action as $a=(i,j,q)$. This names a lossless transfer of the finite quantity $q$ from cell $i$ to cell $j$. The notation is a compact instruction for a represented transformation. Its physical validity depends on the world and constraints already declared.

\section{Seven layers of one event}
The working sequence is:

\begin{center}\small
physical state and proposed action\\$\downarrow$\\
physical feasibility\\$\downarrow$\\
valuation\\$\downarrow$\\
separately declared permission or account policy\\$\downarrow$\\
actor choice\\$\downarrow$\\
execution\\$\downarrow$\\
audit
\end{center}

Each layer answers a distinct question. Feasibility asks whether the represented transformation can occur within physical constraints. Valuation asks what its declared potential difference would be. A separate policy may determine which feasible actions are authorised or affordable. The actor's rule then selects an action or declines to act. Execution changes the physical state. Audit checks what occurred against the declared record.

In a real institution, some legal or safety permissions may be checked earlier as well. The conceptual separation remains: a favourable value is not physical capability or legal permission. An unfavourable value does not decide a person's entitlement to essential service.

\section{What local means}
To evaluate a simple two-cell transfer, the account will need the two stocks, their reference parameters, the action quantity and the applicable process-burden declaration. It need not read an unrelated tank whose potential term is unchanged. Later we prove this cancellation; at present it is a requirement to justify, not a licence to ignore affected effects.

Locality follows dependencies, not mere geographic distance. A long pipe can connect two model cells. A nearby third cell might be irrelevant to one action, while a farther cell could be relevant if the declared potential couples it to the action. The information boundary has to match the physical and mathematical description.

The valuation should not quietly consult a simulated future, a global ranking, a person's wealth or an eventual study outcome. A planning system may use different information for its own task, but that is a separate mechanism. Mixing its forecast into an allegedly present local value changes the claim.

\section{Several actions can occur together}
Suppose two actors both propose to deliver water into Vale from separate tanks. The destination changes once under the combined delivery. The pair therefore needs a coherent joint endpoint and a joint feasibility check. Each action being feasible alone does not ensure that their combination respects a shared capacity.

There is no rule in this introductory world that only one action may be accepted per micro-step. Simultaneous groups remain in scope. We will learn in Chapter~\ref{ch:groups} why simply adding independent values evaluated from the same start can miscount the joint change.

If one delivery truly occurs after another, the second begins from the state the first left. That is a sequential description. It is not equivalent to pretending that both began from the original state. Nor may we invent a physical order merely because a computer happens to list one actor first.

\section{Keep transformation chains visible}
The clinic's real pump belongs to a chain: an energy source supplies electricity, equipment converts it into mechanical work, and a pipe delivers water. When a model represents those processes, their separate physical transformations and boundaries should remain traceable. One purchase need not become one giant physical action that hides all the intermediate effects.

Conversely, several links do not justify charging an extra burden merely because there are several links. An effect must have a physical meaning and be counted once. A regional audit must not count an internal transfer again as new external supply.

\section{An actor rule is an additional choice}
A researcher could specify random choice, a service priority, a planner or a value-maximising policy. Those policies can produce different outcomes even with the same value equation. EBU calculation does not choose among them.

The prospective Gaussian programme considers EBU-blind candidate formation and random choice after separate feasibility and affordability checks. It does not secretly rank by EBU or add a restoring force. This is a proposed behavioural study, not a prescription for patient care or an implemented economy.

\practice{An action has a positive calculated value but would overflow the destination. Should the actor treat the value as permission?}{No. The action fails the physical constraints of the declared world. Valuation cannot make an invalid endpoint feasible. Equally, a feasible negative-valued action raises a separate policy question; its sign alone is not a universal prohibition.}

\carry{We now know what an action and an actor are. To evaluate the action, we need to declare the ruler by which a represented change counts as nearer to or farther from a reference.}
